% =====================================================================
% Exemple minimal — POLYCOPIÉ
% Compilation : make poly (depuis examples/), ou latexmk -pdf main.tex
%                avec TEXINPUTS pointant sur tex/ et assets/.
% =====================================================================
\documentclass[11pt,twoside]{ocots-book}
\usepackage[lang=fr, theme=n7, solutions=end, math=analysis,
            institution=n7,insa, author={Prénom Nom}]{ocots}

\title{Exemple de polycopié}
\subtitle{Sous-titre facultatif}
\author{Prénom \textsc{Nom}}
\date{\today}

\makeindex

\begin{document}

\dominitoc%

% ---------------------------------------------------------------- front
\frontmatter%
\maketitle
\tableofcontents

% ---------------------------------------------------------------- main
\mainmatter%

\part{Première partie}%
\label{part:one}

\ocotscollectsolutions% début de la collecte des corrigés

\chapter{Premier chapitre}%
\label{chap:one}
\minitoc%

\begin{quotation}
    Chapeau du chapitre.
\end{quotation}

\section{Première section}%
\label{sec:one}

Texte, avec une norme $\norm{x}$ et un intervalle $\intervalleoo{a}{b}$.

\begin{theorem}{Titre du théorème}{exemple}
    Énoncé.
\end{theorem}

\begin{proof}
    Démonstration.
\end{proof}

\begin{exercise}[label=ex:demo, points=4]
    Énoncé de l'exercice.
\solution
    Corrigé, reporté en fin de partie par l'option \texttt{solutions=end}.
\end{exercise}

% <<< étape 3 : tous les environnements, toutes les options >>>

%-------------------------------------------
\chapter{Corrections des exercices}
\ocotsprintsolutions% les corrigés collectés viennent ici

% ---------------------------------------------------------------- back
\backmatter%
\printindex%

\end{document}
